% ============================================
% LUC-Escolha (ZH)
% 中文翻译 — SPEC-935-R356
% ============================================
\section*{\textbf{LUC-Escolha} --- 抉择}
\markboth{LUC-Escolha --- 抉择}{}

\noindent \textit{如何处置你随身携带的魔鬼}

\noindent \textbf{2026年4月 --- 巴尔巴塞纳}

\vspace{0.5em}\noindent\rule{\textwidth}{0.5pt}\vspace{0.5em}

\noindent 卢西亚在巴尔巴塞纳唯一的旅馆过了一夜。小房间，墙上有霉斑，天花板上的风扇像缓慢的旋转木马一样转着。

\noindent 她把两本日记放在桌上：原版（棕色封皮，62年的书写）和第二本（黑色封皮，死后写成）。

\noindent 她没有打开第二本。还没有。

\noindent 而是坐在床上，思考。

\noindent 她有三个选项。

\noindent \textbf{选项一：把日记还给档案馆。}

\noindent 假装什么都没发生。让另一个研究者找到它。循环会继续——但她就不必负责了。或者，她负责吗？她已经读过了。污染已经开始。归还日记不会抹去症状。

\noindent \textbf{选项二：销毁日记。}

\noindent 烧掉。撕碎。用酸溶解。如果物理载体被摧毁，循环也许就能停止。但第二本日记呢？如果销毁原版日记释放的是怪物，而不是囚禁它呢？

\noindent \textbf{选项三：接受。}

\noindent 接受循环现在已成为她的一部分。接受右眼正在改变。接受气味不会离去。接受从那一刻起，她是日记的新守护者——而且迟早，她必须把它传下去。

\noindent \textbf{--- 但还有第四个选项。}

\noindent 读第二本日记。不是作为对恐惧的反应——而是作为调查行为。主动选择进入循环，而不是被它拖入。拒绝扮演受害者，承担见证者的角色。卢西亚想，区别就是溺水与潜水的区别。

\vspace{0.5em}\noindent\rule{\textwidth}{0.5pt}\vspace{0.5em}

\noindent 答案在凌晨3点14分到来。

\noindent 气味唤醒了她。

\noindent 不是之前那种甜味。更浓。更稠。湿土的气味。烧过的骨头的气味。东西从地里出来的气味。

\noindent 卢西亚睁开眼睛。

\noindent 第二本日记摊开在桌上。

\noindent 她不记得打开过它。

\noindent 那一页上写着一句话，同样的颤抖笔迹：

\noindent \textit{\textquotedblleft{}它已经替你选择了。你只是还不知道。\textquotedblright{}}

\noindent 卢西亚合上日记。把两本都收进背包。结了旅馆的账。搭上早上5点回贝洛奥里藏特的大巴。

\noindent 整个旅程中，右眼都在灼烧。

\noindent 仿佛有什么东西正在它里面生长。

\noindent 仿佛有什么东西正在\textit{学着看}。

\vspace{0.5em}\noindent\rule{\textwidth}{0.5pt}\vspace{0.5em}

\vspace{0.5em}
\noindent\textit{\textbf{↪ 路线:} \rota{LUC-06} \rotasep \rota{LUC-Entrevista}}

\vspace{0.5em}\noindent\rule{\textwidth}{0.5pt}\vspace{0.5em}

\subsection*{污染量表 —— 更新1}

\noindent \textit{档案已检测到：您跟随卢西亚的调查，一直读到了作出选择的那一刻。}

\begin{moltabletwo}
\hline
分数 & 诊断 \\
\hline
0--1 & 读者完好。不太可能，但仍有可能。 \\
\hline
2--4 & 读者已暴露。您已经读了本不该读的文件。 \\
\hline
5--7 & 读者已被污染。卢西亚作了选择。您也已经作了您的选择——即便您不知道那是哪一个。 \\
\hline
8--10 & 患者1,263正在形成。那生物已经知道您的名字。 \\
\hline
\end{moltabletwo}

\begin{quote}
档案注：分数不累加。量表会随着档案在您周围合拢而自行调整。
\end{quote}
